\documentclass[aps,pre,reprint,superscriptaddress]{revtex4-2}

\usepackage{graphicx}
\usepackage{amsmath}
\usepackage{booktabs}
\usepackage{hyperref}

\begin{document}

\title{Topological Link Prediction for Metabolic Gap-Filling}

\author{Dylan Smith}
\affiliation{Indiana University Bloomington}
\date{\today}

\begin{abstract}
Genome-scale metabolic models are powerful tools for metabolic engineering, but they remain
incomplete because many enzymatic functions are missing or poorly annotated. This project evaluates
whether graph topology can recover hidden metabolic relationships in the \textit{Escherichia coli}
K-12 iJO1366 model. We represent the model as a bipartite metabolite-reaction graph, hide known
metabolite-reaction edges, and compare Node2Vec link prediction against simpler topology-based
baselines. The results are used to assess whether network structure provides useful gap-filling
signal independent of sequence homology.
\end{abstract}

\maketitle

\section{Introduction}

Genome-scale metabolic models encode the reactions, metabolites, and gene associations known for an
organism. Even curated models such as iJO1366 remain incomplete, which limits simulation accuracy
and creates uncertainty in downstream metabolic engineering work. This project asks whether network
topology can provide an independent signal for identifying missing metabolic structure.

\section{Data}

The primary data source is the iJO1366 reconstruction of \textit{E. coli} K-12 MG1655 from BiGG
Models. The model contains reactions, metabolites, stoichiometric coefficients, reversibility, and
gene-protein-reaction rules. Optional validation data include RefSeq NC\_000913.3 and enzyme
annotations from KEGG or BioCyc/EcoCyc.

\section{Methods}

We convert iJO1366 into a bipartite graph with metabolite nodes and reaction nodes. An edge connects
a metabolite to a reaction when the metabolite participates in that reaction. The main benchmark
hides a fraction of known metabolite-reaction edges while keeping all evaluated nodes present in the
training graph.

Baseline scores are computed using simple NetworkX-derived topology features. Node2Vec embeddings
are trained on the training graph, and candidate metabolite-reaction pairs are scored using a
supervised classifier trained on positive and negative examples.

\section{Results}

This section will report graph summary statistics, degree distributions, precision-recall curves,
AUPRC, precision at selected ranks, and top-ranked candidate predictions.

\section{Discussion}

The discussion will interpret whether topology-based link prediction performs better than random and
degree-based baselines, where it succeeds, where it fails, and what the results imply for metabolic
gap-filling.

\section{Limitations}

Important limitations include the use of simulated missingness, possible bias from high-degree
currency metabolites, and the difference between predicting graph structure and proving biochemical
activity.

\section{Code Availability}

Code and reproducible scripts are available at: \url{https://github.com/dylan33smith/TODO}

\bibliography{references}

\end{document}
